\documentclass[12pt]{beamer}
\usetheme{Boadilla}
\usecolortheme{default}

\usepackage[brazil]{babel}
\usepackage[utf8]{inputenc}
\usepackage{amsmath, amssymb}
\usepackage{tikz}
\usetikzlibrary{arrows.meta, calc}

% ── Paleta de cores ──────────────────────────────────────────────
\definecolor{azulEDO}{RGB}{30, 80, 150}
\definecolor{verdeHip}{RGB}{20, 120, 80}
\definecolor{vermelhoAlerta}{RGB}{180, 40, 40}
\definecolor{cinzaFundo}{RGB}{245, 245, 248}

\setbeamercolor{frametitle}{bg=azulEDO, fg=white}
\setbeamercolor{title}{fg=azulEDO}
\setbeamercolor{structure}{fg=azulEDO}
\setbeamercolor{block title example}{bg=verdeHip!80!black, fg=white}
\setbeamercolor{block body example}{bg=verdeHip!10}
\setbeamercolor{block title alerted}{bg=vermelhoAlerta!85, fg=white}
\setbeamercolor{block body alerted}{bg=vermelhoAlerta!8}

% ── Metadados ────────────────────────────────────────────────────
\title{Introdução às Equações Diferenciais}
\subtitle{Aula 1 — O que é uma equação diferencial e por que ela aparece}
\author{}
\institute{Equações Diferenciais Ordinárias}
\date{}

% ── Macro para hipótese ──────────────────────────────────────────
\newcommand{\hipotese}[4]{%
  \begin{block}{Hipótese #1 — #2}
    \textbf{Enunciado:} #3\\[4pt]
    \textbf{Custo:} #4
  \end{block}%
}

% ── Suprime botões de navegação ──────────────────────────────────
\setbeamertemplate{navigation symbols}{}
\setbeamertemplate{footline}[frame number]

% ─────────────────────────────────────────────────────────────────
\begin{document}

\begin{frame}
  \titlepage
\end{frame}

% ================================================================
%  BLOCO 1 — PROBLEMA CONCRETO
% ================================================================

\begin{frame}{Uma pergunta simples — sem equações}

  \begin{center}
    {\large Você prepara um café a \textbf{90\,°C}.}\\[6pt]
    A sala está a \textbf{20\,°C}.\\[14pt]
    \textit{Quanto tempo leva até o café atingir \textbf{60\,°C}?}
  \end{center}

  \vfill
  \pause

  \begin{columns}[T]
    \column{0.48\textwidth}
      \textbf{O que você sabe:}
      \begin{itemize}
        \item O café começa quente e esfria.
        \item Quanto maior a diferença de temperatura, mais rápido esfria.
        \item Eventualmente chega à temperatura da sala.
      \end{itemize}

    \column{0.48\textwidth}
      \textbf{O que você \emph{não} sabe:}
      \begin{itemize}
        \item Em que instante exatamente?
        \item A que taxa o resfriamento ocorre?
        \item Como calcular isso?
      \end{itemize}
  \end{columns}

  \vfill
  \begin{alertblock}{}
    \centering
    Para responder com precisão, precisamos de uma linguagem melhor.
  \end{alertblock}

\end{frame}

% ----------------------------------------------------------------
\begin{frame}{O que a intuição diz sobre a taxa de resfriamento?}

  Seja $T(t)$ a temperatura do café no instante $t$ (em minutos).

  \vspace{10pt}
  A linguagem verbal já contém a ideia-chave:

  \vspace{8pt}
  \begin{exampleblock}{}
    \centering
    \textit{``quanto maior a diferença, mais rápido esfria''}
  \end{exampleblock}

  \vspace{10pt}
  \pause

  Em símbolos, isso é uma afirmação sobre a \textbf{taxa de variação} de $T$:
  \[
    \frac{dT}{dt} \quad \text{depende de} \quad T(t) - T_m
  \]
  onde $T_m = 20\,°C$ é a temperatura do meio.

  \vspace{10pt}
  \pause

  \begin{alertblock}{Pergunta}
    Qual é exatamente essa dependência?
    \quad Como transformar a intuição em equação?
  \end{alertblock}

\end{frame}

% ================================================================
%  BLOCO 2 — HIPÓTESES E EMERGÊNCIA DA EDO
% ================================================================

\begin{frame}{Hipótese 1 — Temperatura uniforme}

  \hipotese{1}{Temperatura uniforme}{%
    O café tem \textbf{uma única temperatura} em cada instante $t$:
    a função $T(t)$ está bem definida.%
  }{%
    Não modelamos gradientes internos nem a xícara em si.%
  }

  \vspace{10pt}
  \pause

  \textbf{Ganho:} $T$ é uma função de \emph{uma} variável — $t$.

  \vspace{8pt}
  \textbf{Quando é razoável:}
  \begin{itemize}
    \item Xícara pequena, líquido bem misturado.
    \item Não vale para um forno ou para o oceano.
  \end{itemize}

  \vspace{8pt}
  \textbf{Consequência matemática:} podemos falar de $\dfrac{dT}{dt}$.

\end{frame}

% ----------------------------------------------------------------
\begin{frame}{Hipótese 2 — Lei de Newton do Resfriamento}

  \hipotese{2}{Proporcionalidade}{%
    A taxa de variação de $T$ é \textbf{proporcional}
    à diferença $T(t) - T_m$.%
  }{%
    Falha para grandes diferenças de temperatura
    ou se $T_m$ variar ao longo do tempo.%
  }

  \vspace{10pt}
  \pause

  \textbf{Tradução imediata:}
  \[
    \frac{dT}{dt} = k\bigl(T(t) - T_m\bigr), \qquad k < 0.
  \]

  \vspace{6pt}
  \pause

  \begin{exampleblock}{A equação emergiu — não foi postulada}
    Cada símbolo carrega uma hipótese:
    \begin{itemize}
      \item $T(t)$: Hipótese 1 (temperatura bem definida).
      \item $k$: Hipótese 2 (proporcionalidade; $k<0$ porque o café esfria).
      \item $T_m$: temperatura do meio, considerada constante.
    \end{itemize}
  \end{exampleblock}

\end{frame}

% ================================================================
%  BLOCO 3 — O NOVO OBJETO MATEMÁTICO
% ================================================================

\begin{frame}{Compare: dois tipos de equação}

  \begin{columns}[T]

    \column{0.45\textwidth}
      \begin{block}{Equação algébrica}
        \[
          x^2 + 5x + 4 = 0
        \]
        \vspace{4pt}
        \textbf{Incógnita:} um número $x$.\\[4pt]
        \textbf{Solução:} $x = -1$ ou $x = -4$.
      \end{block}

    \column{0.45\textwidth}
      \begin{block}{Equação diferencial}
        \[
          \frac{dT}{dt} = k(T - T_m)
        \]
        \vspace{4pt}
        \textbf{Incógnita:} uma \emph{função} $T(t)$.\\[4pt]
        \textbf{Solução:} uma função que satisfaz a relação.
      \end{block}

  \end{columns}

  \vspace{14pt}
  \pause

  \begin{exampleblock}{Conclusão 1}
    Uma \textbf{equação diferencial (ED)} é uma equação que envolve
    uma função desconhecida e suas derivadas.
    A incógnita é uma \emph{função}, não um número.
  \end{exampleblock}

\end{frame}

% ----------------------------------------------------------------
\begin{frame}{Classificação: vocabulário para descrever EDOs}

  Vamos falar sobre a equação que já temos:
  \[
    \frac{dT}{dt} = k(T - T_m).
  \]

  \vspace{6pt}
  \begin{columns}[T]

    \column{0.48\textwidth}
      \textbf{Tipo}\\[4pt]
      Derivadas em relação a \emph{uma} variável ($t$)?
      \[
        \Rightarrow \text{EDO (ordinária)}
      \]
      \vspace{4pt}
      Derivadas parciais?
      \[
        \Rightarrow \text{EDP (parcial)}
      \]

    \column{0.48\textwidth}
      \textbf{Ordem}\\[4pt]
      Maior derivada presente: $\dfrac{dT}{dt}$ (1ª derivada).
      \[
        \Rightarrow \text{Ordem } \mathbf{1}
      \]

      \vspace{6pt}
      \textbf{Linearidade}\\[4pt]
      $T$ e $T'$ aparecem ao 1º grau, coeficientes dependem de $t$.
      \[
        \Rightarrow \text{Linear}
      \]

  \end{columns}

\end{frame}

% ----------------------------------------------------------------
\begin{frame}{Mais exemplos de classificação}

  \begin{center}
  \renewcommand{\arraystretch}{2.0}
  \begin{tabular}{lcc}
    \hline
    \textbf{Equação} & \textbf{Ordem} & \textbf{Linear?} \\
    \hline
    $\dfrac{dT}{dt} = k(T - T_m)$
      & 1 & Sim \\
    $\dfrac{d^2y}{dx^2} - 2\dfrac{dy}{dx} + y = 0$
      & 2 & Sim \\
    $\dfrac{dy}{dx} = xy^{1/2}$
      & 1 & Não \\
    $\dfrac{d^2y}{dx^2} + 5\!\left(\dfrac{dy}{dx}\right)^{\!3} - 4y = e^x$
      & 2 & Não \\
    \hline
  \end{tabular}
  \end{center}

  \vspace{8pt}
  \pause

  \begin{alertblock}{Regra para linearidade}
    A EDO de ordem $n$ é \textbf{linear} se $y, y', \ldots, y^{(n)}$
    aparecem ao \textbf{1º grau} e os coeficientes dependem
    \emph{apenas} da variável independente.
  \end{alertblock}

\end{frame}

% ================================================================
%  BLOCO 4 — SOLUÇÃO E PVI
% ================================================================

\begin{frame}{O que é uma solução?}

  De volta ao café.
  \[
    \frac{dT}{dt} = k(T - 20).
  \]

  \pause
  \vspace{6pt}
  \textbf{Afirmação:} a função $T(t) = 20 + 70\,e^{kt}$ satisfaz a equação.

  \vspace{6pt}
  \textbf{Verifique:}
  \[
    \frac{dT}{dt} = 70k\,e^{kt}, \qquad
    k(T - 20) = k \cdot 70\,e^{kt}. \quad \checkmark
  \]

  \pause
  \vspace{6pt}
  \begin{exampleblock}{Definição — Solução de uma EDO}
    Uma função $\phi(t)$, definida em um intervalo $I$,
    é \textbf{solução} da EDO se, ao substituir $\phi$ na equação,
    ela se torna uma \textbf{identidade} em $I$.
  \end{exampleblock}

  \pause
  \vspace{6pt}
  Note que $T(t) = 20 + Ce^{kt}$ é solução para \emph{qualquer} constante $C$.
  Isso é a \textbf{família de soluções (solução geral)}.

\end{frame}

% ----------------------------------------------------------------
\begin{frame}{Problema de Valor Inicial (PVI)}

  A família $T(t) = 20 + Ce^{kt}$ tem infinitas soluções.

  \vspace{8pt}
  \textbf{A pergunta original exige mais:}
  o café começa a \textbf{90\,°C} no instante $t = 0$.

  \[
    \frac{dT}{dt} = k(T - 20), \qquad \underbrace{T(0) = 90}_{\text{condição inicial}}.
  \]

  \pause
  \vspace{6pt}
  Aplicando: $T(0) = 20 + C = 90 \;\Rightarrow\; C = 70$.

  \vspace{6pt}
  \[
    \boxed{T(t) = 20 + 70\,e^{kt}}
  \]

  \pause
  \vspace{6pt}
  \begin{exampleblock}{Definição — PVI}
    Um \textbf{problema de valor inicial} é uma EDO acompanhada de
    condições sobre a função (e suas derivadas) em um único ponto $t_0$.
    Resolver o PVI é encontrar a \textbf{solução particular}.
  \end{exampleblock}

\end{frame}

% ================================================================
%  BLOCO 5 — VISUALIZAÇÃO QUALITATIVA
% ================================================================

\begin{frame}{O que a equação já nos diz — sem resolver}

  \[
    \frac{dT}{dt} = k(T - 20), \qquad k < 0.
  \]

  \vspace{6pt}
  Análise de sinal sem integrar:

  \vspace{8pt}
  \begin{columns}[T]

    \column{0.48\textwidth}
      \textbf{Se $T > 20$:}\\
      $T - 20 > 0$ e $k < 0$\\
      $\Rightarrow \dfrac{dT}{dt} < 0$\\[6pt]
      O café \textbf{esfria}. ✓

    \column{0.48\textwidth}
      \textbf{Se $T < 20$:}\\
      $T - 20 < 0$ e $k < 0$\\
      $\Rightarrow \dfrac{dT}{dt} > 0$\\[6pt]
      Um objeto frio \textbf{aquece}. ✓

  \end{columns}

  \vspace{10pt}
  \pause

  \begin{exampleblock}{Equilíbrio}
    $\dfrac{dT}{dt} = 0 \;\Leftrightarrow\; T = 20$. \\[4pt]
    A solução constante $T(t) = 20$ é chamada de
    \textbf{solução de equilíbrio}.
  \end{exampleblock}

  \vspace{4pt}
  \pause
  \small
  Essa análise qualitativa — sem resolver — é uma das três abordagens
  do curso: \textbf{analítica}, \textbf{qualitativa}, \textbf{numérica}.

\end{frame}

% ================================================================
%  BLOCO 6 — O MESMO MODELO, FENÔMENOS DIFERENTES
% ================================================================

\begin{frame}{Uma equação, muitos fenômenos}

  \[
    \frac{dy}{dt} = ky
  \]

  \vspace{6pt}
  \begin{columns}[T]

    \column{0.30\textwidth}
      \textbf{Crescimento\\populacional}\\[4pt]
      $y = P(t)$: população.\\[4pt]
      $k > 0$: crescimento.\\[4pt]
      \textit{(Malthus, 1798)}

    \column{0.30\textwidth}
      \textbf{Decaimento\\radioativo}\\[4pt]
      $y = A(t)$: massa.\\[4pt]
      $k < 0$: decaimento.\\[4pt]
      \textit{(datação C-14)}

    \column{0.30\textwidth}
      \textbf{Juros\\compostos}\\[4pt]
      $y = S(t)$: capital.\\[4pt]
      $k = r > 0$: taxa.\\[4pt]
      \textit{(capitalização)}

  \end{columns}

  \vspace{14pt}
  \pause

  \begin{exampleblock}{Conclusão 2}
    Uma única EDO pode servir como modelo matemático
    para fenômenos muito diferentes.
    O que muda é a \textbf{interpretação dos símbolos}.
  \end{exampleblock}

\end{frame}

% ================================================================
%  BLOCO 7 — SÍNTESE
% ================================================================

\begin{frame}{Síntese — Como chegamos até aqui}

  \begin{enumerate}
    \setlength\itemsep{8pt}
    \item Partimos de um \textbf{problema concreto}: resfriamento do café.
    \item Formulamos \textbf{duas hipóteses} (temperatura uniforme;
          resfriamento proporcional à diferença).
    \item A EDO \textbf{emergiu} das hipóteses:
          $\dfrac{dT}{dt} = k(T - T_m)$.
    \item Definimos \textbf{solução} como função que satisfaz
          a equação identicamente.
    \item O \textbf{PVI} seleciona uma solução particular da família.
    \item A \textbf{análise qualitativa} extraiu informação
          \emph{antes} de resolver.
  \end{enumerate}

  \vspace{10pt}
  \pause

  \begin{columns}[T]
    \column{0.48\textwidth}
      \textbf{Vocabulário da aula:}
      \begin{itemize}
        \item EDO / EDP, ordem, linearidade
        \item Solução geral, solução particular
        \item Problema de valor inicial (PVI)
        \item Solução de equilíbrio
      \end{itemize}
    \column{0.48\textwidth}
      \textbf{Três abordagens do curso:}
      \begin{itemize}
        \item Analítica (resolver explicitamente)
        \item Qualitativa (análise de sinal, campo de direções)
        \item Numérica (aproximar a solução)
      \end{itemize}
  \end{columns}

\end{frame}

% ================================================================
%  BLOCO 8 — PONTE PARA A PRÓXIMA AULA
% ================================================================

\begin{frame}{Próxima aula — Uma pergunta que ainda não respondemos}

  Temos a solução $T(t) = 20 + 70\,e^{kt}$, mas ainda não sabemos $k$.

  \vspace{10pt}
  Para determinar $k$ a partir de dados, precisamos \textbf{resolver} a equação
  de forma sistemática.

  \vspace{10pt}
  \pause

  \begin{alertblock}{Pergunta em aberto}
    $\dfrac{dT}{dt} = k(T - 20)$ pode ser reescrita como
    \[
      \frac{dT}{T - 20} = k\,dt.
    \]
    \vspace{4pt}
    Isso é uma coincidência útil ou existe uma classe inteira de
    equações que admitem essa separação?
  \end{alertblock}

  \vspace{10pt}
  \pause

  \begin{center}
    \textbf{Próxima aula:} Equações Separáveis.\\[4pt]
    \small Pré-requisito: técnicas de integração (substitução, frações parciais).
  \end{center}

\end{frame}

% ================================================================
%  APÊNDICE — EXERCÍCIOS
% ================================================================

\begin{frame}{Exercícios}

  \begin{enumerate}
    \setlength\itemsep{8pt}

    \item Verifique que $y = e^{0{,}1x^2}$ é solução de
          $\dfrac{dy}{dx} = 0{,}2xy$ no intervalo $(-\infty, +\infty)$.

    \item Classifique as equações abaixo quanto a tipo, ordem e linearidade:
          \begin{enumerate}[(a)]
            \item $\dfrac{d^2x}{dt^2} + 16x = 0$
            \item $\dfrac{\partial^2 u}{\partial x^2}
                  + \dfrac{\partial^2 u}{\partial y^2} = 0$
            \item $(1-y)\,y' + 2y = e^x$
          \end{enumerate}

    \item Mostre que $y = xe^x$ é solução de $y'' - 2y' + y = 0$.

    \item Para o modelo $\dfrac{dP}{dt} = 0{,}15P + 20$, $P(0) = 100$:
          \begin{enumerate}[(a)]
            \item Qual é a taxa de crescimento em $t = 0$?
            \item Sem resolver, diga se $P(t)$ é crescente ou decrescente.
          \end{enumerate}

  \end{enumerate}

\end{frame}

\end{document}
